% Đặc tuyến lượng tử hóa tuyến tính và sai số lượng tử hóa.
\begin{tikzpicture}[
    x=1.45cm,
    y=1.25cm,
    font=\large,
    line cap=round,
    line join=round
]
    \begin{scope}[draw=black, text=black]
        % Hai trục tọa độ, không có đầu mũi tên.
        \draw[line width=1.2pt] (-4.45, 0) -- (2.90, 0);
        \draw[line width=1.2pt] (0, -0.50) -- (0, 3.90);

        % Đặc tuyến bậc thang với các bước lượng tử bằng nhau.
        \draw[line width=1.5pt]
            (-2.75, 0) -- (-2.75, 1) -- (-1.75, 1)
            -- (-1.75, 2) -- (-0.75, 2)
            -- (-0.75, 3) -- (0.25, 3)
            -- (0.25, 4) -- (1.30, 4);

        % Độ cao một mức lượng tử.
        \draw[<->, >=stealth, line width=1.05pt]
            (-1.80, 0.025) -- (-1.80, 0.975);
        \node[anchor=west, inner sep=0pt]
            at (-1.48, 0.48) {$\Delta x$};

        % Giá trị đầu vào nằm giữa hai mức và sai lệch so với mức dưới.
        \draw[fill=white, line width=0.9pt]
            (0, 3.35) circle[radius=2.4pt];
        \draw[line width=0.8pt] (-0.13, 3.35) -- (0.13, 3.35);
        \draw[line width=0.8pt] (0.38, 3.35) -- (0.65, 3.35);
        \draw[line width=0.8pt] (0.38, 3.00) -- (0.65, 3.00);
        \draw[<->, >=stealth, line width=0.6pt]
            (0.49, 3.025) -- (0.49, 3.325);
        \node[anchor=west, inner sep=0pt]
            at (0.73, 3.16) {$\delta x$};

        % Mũi tên dẫn tới chú thích sai số lượng tử hóa.
        \draw[->, >=stealth, line width=0.95pt]
            (0.89, 2.98) .. controls (1.05, 2.80) and (1.13, 2.52)
            .. (1.16, 2.30);
        % Căn giữa chú thích và công thức trong cùng một khối.
        \node[anchor=north, align=center, inner sep=0pt, font=\large\bfseries]
            at (2.30, 2.15) {Sai số lượng tử hóa\\[6pt]
                {\normalfont\Large $\leq \dfrac{\Delta x}{2}$}};
    \end{scope}
\end{tikzpicture}
